\section{Primera fase: entender y preparar los datos}\label{sec:datos}

\begin{idea}{Entrada, trabajo y salida}
\textbf{Entran} siete CSV del profesor. \textbf{Se comprueban} claves, periodos,
ausencias y balances; se separan clientes y automatriculaciones y se estudian los
episodios excepcionales. \textbf{Salen} una serie mensual por segmento, un
diagnóstico de A012 y alternativas explícitas para lo que no se observa.
\end{idea}

\subsection{Qué contiene cada fuente}
Los ficheros no empiezan todos en el mismo año. Mercado y economía comienzan en
A001; las ventas internas, en A003; el stock, en A009. Una fila de ventas tiene
tres claves: mes, segmento y canal. Esto explica las \num{2520} filas:
$10\text{ años}\times12\text{ meses}\times7\text{ segmentos}\times3\text{ canales}$.
La tabla de stock no tiene canal, de modo que no se puede cruzar con ventas como
si cada canal dispusiera de su propio inventario.

\tabla{inventario.tex}{Datos originales disponibles y granularidad de cada fila.}
{tab:inventario}{\archivo{data/README.md} y anexo 5 del enunciado}

Los años A001--A013 son códigos del caso, no años naturales identificables. El
contador $t$ permite ordenar los meses: diciembre de A012 es $t=144$ y enero de
A013 es $t=145$. Las fechas sintéticas que utiliza una biblioteca no proporcionan
el calendario laboral: días laborables y Semana Santa se leen en los CSV.

\subsection{Por qué limpiar no significa cambiar los datos}
La lectura preserva el código de modelo \archivo{NA}; no lo convierte en un
valor ausente. Un campo vacío sí representa información no disponible. Un cero,
en cambio, es una cantidad observada: por ejemplo, ventas nulas antes de un
lanzamiento. Las columnas porcentuales tampoco comparten siempre escala:
\archivo{pct_diesel} se expresa como proporción, mientras el paro se expresa en
puntos de porcentaje.

Las comprobaciones verifican que las claves no estén duplicadas, que los periodos
sean coherentes y que las sumas y desgloses tengan sentido. Pasan los 25 controles
de lectura y consistencia. Esto \textbf{no implica que todas las identidades
históricas sean exactas}: las discrepancias se muestran sin alterar los originales.

\paso{El balance histórico.}
Si $S_t$ es el stock al cerrar el mes, $L_t$ la facturación a la red y $M_t$ las
matriculaciones, el residuo es
\begin{equation}
 r_t=S_t-\bigl(S_{t-1}+L_t-M_t\bigr).
\end{equation}
En el histórico existen diferencias pequeñas y algunos casos mayores: el máximo
absoluto es de 51 coches en SSUV y de 8 en MD. En MSUV el máximo es 5; en CSUV,
cero. Se mantienen las discrepancias, con localización en las tablas de diagnóstico.
No se fuerza el residuo a cero para mejorar el aspecto de los datos.

\fuente{\archivo{resultados/datos/tablas/balance_stock_segmento.csv}}

Es importante usar aquí matriculaciones y no sustituirlas por ventas netas. Las
automatriculaciones forman parte de las salidas registradas en el balance
histórico. Si se retiran sólo de las salidas y no se cambia el sistema de
contabilidad del stock, se introduce un descuadre artificial.

\subsection{De matriculaciones a compras observadas de clientes}
Para cada mes y segmento se define
\begin{equation}\label{eq:neto}
 O_{s,t}=M_{s,t}-A_{s,t},
\end{equation}
donde $O$ son las compras observadas de clientes, $M$ las matriculaciones y $A$
las automatriculaciones; todas las cantidades se expresan en coches. Las ventas
especiales no se restan automáticamente: una venta a una empresa puede ser una
compra real de un cliente.

En A012 se registran \num{23175} matriculaciones y \num{2469}
automatriculaciones. Por tanto, se observan \num{20706} compras de clientes.
En A011 se observan \num{19543}. La variación es
\begin{equation}
 g_{\mathrm{observado}}=
 \frac{\num{20706}}{\num{19543}}-1\simeq 6{,}0\,\%.
\end{equation}
La serie resultante sigue sin ser toda la demanda potencial. Puede incluir
compras adelantadas y puede estar limitada por falta de coches.

\subsection{El lanzamiento de CSUV cambia la comparación}
CSUV empieza a venderse en agosto de A012 y aporta \num{1192} clientes ese año.
Son clientes reales y pertenecen al total de Kaido. Al excluirlos queda una gama
comparable con \num{19514} clientes, prácticamente estable frente a A011:
$-0{,}1\,\%$.

No hay contradicción entre que crezca la marca y no crezca su gama anterior. Son
dos preguntas distintas: cuánto vende el conjunto y cómo evoluciona el perímetro
de productos que ya existía.

\tabla{crecimiento.tex}{Crecimiento observado y ajuste parcial de episodios.
Los valores ajustados son estimaciones, no recuentos de clientes identificados.}
{tab:crecimiento}{\archivo{resultados/datos/tablas/crecimiento_demanda_a012.csv}}

\grafica{crecimiento_demanda_observado_ajustado_a012.png}
{A012 frente a A011: la conclusión depende de mantener el perímetro de comparación.
El total incluye CSUV; el ajuste central trata episodios, pero no identifica toda
la demanda potencial.}
{fig:crecimiento}{\archivo{datos.ipynb}, gráfico de crecimiento observado y ajustado}

\subsection{Cómo se identifican las compras adelantadas}
Una compra adelantada cambia de fecha; no crea necesariamente un cliente nuevo.
Los indicadores del caso señalan un episodio en agosto de A011, seguido de caída
en septiembre--diciembre, y otro en noviembre--diciembre de A012.

Para estimar el exceso se utilizan A009 y A010 como referencias estacionales.
Cada referencia se escala \emph{por separado} al nivel observado antes del
episodio del año que se estudia. Para una referencia $r$, un segmento $s$ y un
año $a$, la cuenta conceptual es
\begin{align}
 k_{s,a,r}&=\frac{\text{clientes antes del episodio en }a}
                  {\text{clientes en los mismos meses de }r},\\
 E_{s,a,r}&=\max\!\left(0,\ \text{clientes en meses del episodio de }a
                 -k_{s,a,r}\,\text{clientes en esos meses de }r\right).
\end{align}
El mínimo y máximo proceden de esas dos referencias y el central es su punto
medio. El recorte a cero evita atribuir un adelanto negativo; también recuerda que
la estimación es una regla de identificación, no una observación directa.

\tabla{adelantos.tex}{Exceso estimado en los episodios y compensación posterior.
La amplitud es sensibilidad a referencias, no un intervalo estadístico.}
{tab:adelantos}{\archivo{resultados/datos/tablas/adelantos_a011_a012.csv}}

\paso{A011: estudiar el pico y la caída conjuntamente.}
Se estiman \num{450.4} coches de exceso y \num{257.3} de caída posterior
compensatoria, limitada al exceso del segmento. La estimación anual ajustada es
\begin{equation}
 \num{19543}-\num{450.4}+\num{257.3}\simeq\num{19349.9}.
\end{equation}
El saldo no compensado no demuestra que esas compras se tomaran de A012: el
episodio fue en agosto y la referencia estacional es imperfecta. Por eso no se
convierte automáticamente en una transferencia cierta al año siguiente.

\paso{A012: registrar una sola devolución hacia A013.}
El exceso central de cierre se estima en \num{878.4} coches. Al restarlo una
sola vez se obtienen \num{19827.6} compras ajustadas en A012. Comparadas con
las \num{19349.9} ajustadas de A011, implican un aumento parcial del 2,5\,\%.

El registro mensual reparte el ajuste anual entre los meses del episodio. No
resta el exceso anual completo en cada mes. La misma magnitud alimentará después
la caída supuesta en A013: retirar la anomalía del histórico y representar su
devolución futura cumplen funciones diferentes, pero no deben generar dos
descuentos sobre la misma previsión futura.

\paso{Por qué CSUV necesita otro tratamiento.}
No existe una historia de ventas propias anterior a su lanzamiento. Se utiliza
como proxy el exceso de su mercado en noviembre--diciembre respecto a referencias
escaladas, aplicado a las compras de clientes del segmento. En esta ejecución no
se detecta un exceso positivo con ese método. Es un resultado del proxy, no una
demostración de que ningún comprador adelantara su compra.

\subsection{Qué sabemos y qué no sabemos sobre la falta de stock}
Una señal de stock libre bajo identifica un mes sospechoso, pero no cuenta a
los clientes que se marcharon. Si se disponía de diez coches y se vendieron
diez, la demanda pudo ser diez o veinte. El volumen perdido no se observa en
esas ventas.

El indicador de preparación compara el stock sin cliente con las
matriculaciones del mes y utiliza medio mes de cobertura como umbral. No es el
nivel de servicio futuro ni el stock de seguridad: sirve para señalar un posible
problema de observación del histórico.

Como sensibilidad se utiliza la relación entre mercado y compras de clientes
en hasta seis meses anteriores sin señal de restricción, adelanto ni caída.
Si esa referencia sugiere una demanda mayor, se construye
\begin{equation}
 O^{\mathrm{alternativa}}_{s,t}
 =\max\!\left(O_{s,t},\frac{I_{s,t}}
 {\operatorname{mediana}_{u<t}\bigl(I_{s,u}/O_{s,u}\bigr)}\right),
\end{equation}
con $I$ igual a las matriculaciones del mercado y usando sólo antecedentes
válidos. Si no hay referencia suficiente, se conserva lo observado.

La alternativa arroja un crecimiento observado ampliado del 6,8\,\% para el
total, frente al 6,0\,\% de las compras registradas. No es una estimación
independiente de la demanda perdida: depende de la estabilidad de esa relación
comercial. La previsión principal se valida contra clientes \emph{observados},
no contra una serie inventada por el propio modelo.

\grafica{stock_red_antiguedad_a011_a012.png}
{Stock observado de la red por edad. La cantidad total y la antigüedad son dos
dimensiones distintas: reducir coches no garantiza que desaparezca primero el
inventario problemático.}
{fig:stock-historico}{\archivo{datos.ipynb}; stock mensual A011--A012}

\begin{idea}{Qué concluye realmente la primera fase}
Se observa crecimiento de clientes en el total y estabilidad de la gama anterior.
El ajuste central de episodios sigue siendo positivo para el total, pero depende
de las referencias y de la compensación de A011. La señal de stock y la historia
breve de CSUV impiden presentar una cifra única como demanda potencial exacta.
La política futura de cero automatriculaciones es un supuesto de negocio, no una
corrección retroactiva de lo ocurrido.
\end{idea}
\FloatBarrier
\grafosanchos{demanda_observada_segmento_a003_a012.png}
{Compras observadas por segmento, A003--A012. Los lanzamientos y los episodios
marcados impiden tratar todas las series como una historia homogénea. Las escalas
son distintas para ver los patrones; no deben compararse alturas entre paneles.}
{fig:historico-clientes}{\archivo{datos.ipynb}; serie mensual neta de automatriculaciones}
